\subsection{Definition of a principal ideal domain}

Recall (I, p.~104) that an ideal of a commutative ring A is said to be principal if it has the form \((a) = A a\) for some \(a \in A\) .

DEFINITION 1. --- A principal ideal domain is an integral domain (I, p.~116) in which every ideal is principal.

Examples. --- The ring Z of rational integers is a principal ideal domain (I, p.~111). If K is a commutative field, then the polynomial ring K{[}X{]} in one indeterminate over K is a principal ideal domain (IV, p.~11, Prop. 11); the same is true of the ring of formal power series K{[}{[}X{]}{]}, for every ideal of this ring has the form (X \(^{n}\) ) (IV, p.~38, Prop. 12). * The ring of integers in a p-adic field is a principal ideal domain. *

If \(\mathbf{Q}(i)\) denotes the field obtained from the field Q of rational numbers by adjoining a root i of the irreducible polynomial \(X^{2}+1\) , then the elements \(a+bi\) of \(\mathbf{Q}(i)\) , where a and b are rational integers, form a subring A of \(\mathbf{Q}(i)\) , called « the Gaussian integers », which is a principal ideal domain (VII, p.~50, Ex. 7). By contrast, in the field \(\mathbf{Q}(\rho)\) , where \(\rho\) is a root of \(X^{2}+5\) , the subring B consisting of elements \(a+b\rho\) (a and b rational integers) is not a principal ideal domain (VII, p.~51, Ex. 12).

The polynomial ring K{[}X, Y{]} in two indeterminates over a field K is not a principal ideal domain. Indeed the nonzero constants are the only elements which divide both X and Y, and none of them generates the ideal generated by X and Y.

